\section{引言}

全切片图像（whole-slide images，WSIs）以极高分辨率记录完整的组织形态，是计算病理学及辅助诊疗的重要基础。准确的组织分割可服务于肿瘤范围测量、肿瘤微环境分析、感兴趣区域筛选和下游诊断建模。然而，单张 WSI 可包含数十亿像素，并具有复杂边界和显著的尺度变化，逐像素人工标注成本高且难以扩展。因此，如何以有限人工交互获得准确、可适配的组织图谱，仍是 WSI 分析中的关键问题。

现有方法通常将组织分割视为预定义类别集合上的多分类预测，其输出仅限于训练阶段见过的类别。面对新的组织结构、特定细胞组成或研究任务定义的目标，模型往往需要新增像素级标注并重新训练。由于组织定义还会随癌种、队列和临床问题而变化，固定类别模型难以支持可扩展的组织分析。

以分割一切模型（Segment Anything Model，SAM）为代表的可提示模型允许用户通过少量点、边界框或涂画指定目标，无需重新训练特定类别的预测头。但直接应用于 WSI 仍面临三项挑战：组织身份不仅取决于颜色、纹理和边界，还与细胞类型、密度及空间组成有关；同类组织可能分散在多个不连通区域，局部提示难以覆盖整张切片；病理组织同时包含细胞形态、局部结构和全局布局，单一尺度难以兼顾边界精度与语义判别。

因此，提示驱动的 WSI 分割不应只是围绕点击位置扩展局部边界，而应从一次或少数几次交互中推断目标组织概念，并检索整张切片中所有语义一致的区域。我们的核心观察是：低倍外观相似的区域可能具有不同的细胞组成，而同类组织的离散实例通常共享细胞与组织结构模式。为此，我们将组织表示为融合视觉轮廓与细胞证据的可学习区域，并在精细、中间和粗粒度尺度间进行层次化推理。

图~\ref{fig:overview} 展示了 FewClick，一个面向 WSI 少提示组织分割的细胞感知层次化区域提示框架。该框架从对齐的多尺度视觉特征中生成可学习组织区域，并通过细胞到区域注意力注入细胞嵌入、密度和组成信息。稀疏父子层次结构与门控适配器进一步融合局部边界、邻域信息和宏观上下文。交互阶段，两个置换不变编码器分别聚合正、负提示；上下文感知解码器将目标表示与全切片区域匹配，再把区域概率投影至像素空间。由此，FewClick 能够检索空间上不连通但具有一致细胞组成和组织语义的区域。

本文的主要贡献包括四个方面：（1）我们将 WSI 分割表述为由提示定义的二值组织检索问题，而非基于固定输出通道的多分类预测；（2）我们提出可微分的细胞感知组织区域，以统一编码视觉特征、细胞嵌入、细胞密度与细胞组成；（3）我们在空间对齐的精细、中间和粗粒度区域之间设计了稀疏的自底向上与自顶向下交互机制；（4）我们将带符号的提示编码与上下文感知区域解码相结合，使正、负证据均能参与整张切片范围内的推理。受控对比实验表明，FewClick 相较三种可提示分割基线均取得了稳定提升；消融实验进一步验证了可学习区域、细胞证据和多尺度上下文各自的贡献。
